\section{评估指标}
\label{sec:metrics-zh}

\paragraph{主体标记诊断。}
标记准确率是未遮蔽乐谱主体目标中获得最大 logit 的比例。交叉熵在所有显式窗口的同一目标标记上累计，条件前缀位置和重叠上下文位置均被遮蔽。

\paragraph{音级集合指标。}
对于生成音级支集 $P(x)$ 和非空目标集合 $S$，覆盖率为 $|P(x)\cap S|/|S|$，精确率为 $|P(x)\cap S|/|P(x)|$，Jaccard 相似度为 $|P(x)\cap S|/|P(x)\cup S|$。音程向量距离是生成与目标六维音程类向量之间的 $L_1$ 距离。若样本没有音级集合目标，这些指标记为不适用。

\paragraph{序列指标。}
循环序列顺序准确率将生成音级与请求变换序列的循环重复位置进行比较。严格序列变换准确率统计与请求形式完全相同且互不重叠的完整十二音块。聚合完成率为十二个音级中实际出现的比例，仅在序列条件样本上取平均。

\paragraph{节奏指标。}
密度曲线统计每小节的音符起音数。节奏轮廓距离将按请求声部数归一化的生成曲线与解码阶段可见的确定性轮廓模板比较。密度曲线误差则把生成曲线与语料生成时保存的随机目标曲线比较；后者不进入模型条件和候选排序。

\paragraph{姿态指标。}
姿态特征包括起音离散度、音区跨度、按请求声部数归一化的音符密度、平均时值和逐声部休止覆盖率。休止覆盖率根据时间区间并集计算，范围限制在 $[0,1]$。各特征接近度的平均值构成启发式姿态一致性分数。该分数具有可解释性，但不是感知评分。

\paragraph{音域、跨度、声部与记谱检查。}
音域违反率是超出请求乐器音域的有音高事件比例。内容跨度遵循率为实际事件结束时间与请求时长之比，并截断至 1。声部遵循率是包含生成材料的请求声部索引比例。MusicXML 评估分别报告结构解析成功率、请求小节数遵循率和请求声部数遵循率。

\paragraph{适用性与汇总。}
每项汇总指标只在对应目标存在的样本上计算。目标缺失时写为不适用，而不是 0 或 1。正式 CSV 的每一行均含全部要求的指标列和恰好 2,000 个已评估条件。下文所用 13 行结果在提升为规范 CSV 和 LaTeX 表格之前均通过来源与完成状态门控。
